%% kapa: exam=2 round=34 year=2023 period=1 subject=appraisal_practice subject_ko=감정평가실무 seq=2 number=2 points=30
%% source: Q-Net 감정평가사 기출문제 (cst00309) artlSeq=5212168 file=raw/1008/5212168/1교시 23년도 감정평가사 2차시험_감정평가실무.pdf
%% conversion: kordoc 4.24.51 -> tools/gen2.py (auto); review=proofread-gpt-sol; flags=none
\begin{longquestion}{【문제 2】}
감정평가사 甲은 ○○가로주택정비사업의 관리처분계획수립을 위한 종전자산 및 종후자산 감정평가를 의뢰받았다. 관련 법규 및 이론에 의거 제시된 자료를 활용하여 다음 물음에 답하시오. (30점) \par
\medskip\noindent 물음 1) 종전자산을 감정평가하시오. (20점) \par
\medskip\noindent 물음 2) 가로주택정비사업에 따른 조합원 F의 분담금(또는 환급금)을 구하시오. (5점) \par
\medskip\noindent 물음 3) 가로주택정비사업의 관리처분계획 수립을 위한 종전자산 감정평가와 종후자산 감정평가 업무수행 시 유의사항을 비교하여 설명하시오. (5점) \par
< 자료 1 > 공통사항 \par
1. 비교사례 선정 시 선정사유 및 제외사유를 기재 \par
2. 요인비교치는 상승식으로 산출하되, 소수점 셋째자리에서 반올림하여 둘째자리까지 표기 \par
3. 감정평가액은 십만원 단위에서 반올림하여 백만원 단위까지 표기 \par
4. 조합원 F는 종후자산(아파트) 중 101동 5층 502호를 분양받을 예정 \par
5. 비례율 산정 시 소수점 셋째자리에서 반올림하여 둘째자리까지 표기 \par
< 자료 2 > 사업추진 경과 \par
- 2022.09.15: 조합설립인가 \par
- 2022.12.15: 감정평가법인 약정 체결 \par
- 2023.02.15: 건축심의 신청 \par
- 2023.03.15: 건축심의 조건부 의결 \par
- 2023.05.15: 건축심의 결과 통지서 수령 \par
- 2023.06.15: 종전자산 현장조사 실시 \par
- 2023.07.15: 종전자산 현장조사 완료 \par
< 자료 3 > 종전자산 평가목록 \par
1. 종전자산 건축물 현황 \par
{\small\centering
\resizebox{\ifdim\width>\linewidth\linewidth\else\width\fi}{!}{%
\begin{tabular}{|c|c|c|}
\hline
\makecell{소재지} & \makecell{평가동 1-1번지} & \makecell{평가동 1-2번지} \\
\hline
\makecell{건물명} & \makecell{○○빌라 가동} & \makecell{○○빌라 나동} \\
\hline
\makecell{구조} & \makecell{연와조 스라브지붕 3층} & \makecell{연와조 스라브지붕 3층} \\
\hline
\makecell{주용도} & \makecell{다세대주택} & \makecell{다세대주택} \\
\hline
\makecell{사용승인일} & \makecell{1985.01.02} & \makecell{1985.01.07} \\
\hline
\makecell{건축면적/연면적} & \makecell{73.2m² / 219.6m²} & \makecell{74.4m² / 223.2m²} \\
\hline
\makecell{세대수} & \makecell{6} & \makecell{6} \\
\hline
\makecell{향} & \makecell{남향} & \makecell{동향} \\
\hline
\makecell{형태} & \makecell{계단식} & \makecell{계단식} \\
\hline
\end{tabular}}\par}

2. 토지 등 소유자 명부 \par
{\small\centering
\resizebox{\ifdim\width>\linewidth\linewidth\else\width\fi}{!}{%
\begin{tabular}{|c|c|c|c|c|c|c|c|c|c|c|c|}
\hline
\multirow{3}{*}{\makecell{연번}} & \multirow{3}{*}{\makecell{부번}} & \multirow{3}{*}{\makecell{조합원}} & \multicolumn{3}{|c|}{\makecell{물건내역}} & \multicolumn{5}{|c|}{\makecell{권리내역}} & \multirow{3}{*}{\makecell{비고}} \\
\cline{4-11}
 &  &  & \multirow{2}{*}{\makecell{동}} & \multirow{2}{*}{\makecell{층}} & \multirow{2}{*}{\makecell{호}} & \multicolumn{2}{|c|}{\makecell{토지(m²)}} & \multicolumn{3}{|c|}{\makecell{건축물(m²)}} &  \\
\cline{7-11}
 &  &  &  &  &  & \makecell{지목} & \makecell{지분} & \makecell{구조} & \makecell{전유} & \makecell{공용} &  \\
\hline
\makecell{1} & \makecell{1} & \makecell{A} & \makecell{가} & \makecell{1} & \makecell{101} & \makecell{대} & \makecell{24.6} & \makecell{연와조} & \makecell{32.2} & \makecell{4.4} &  \\
\hline
\makecell{2} & \makecell{1} & \makecell{B} & \makecell{가} & \makecell{1} & \makecell{102} & \makecell{대} & \makecell{12.3} & \makecell{연와조} & \makecell{16.1} & \makecell{2.2} & \makecell{공유} \\
\hline
\makecell{2} & \makecell{2} & \makecell{C} & \makecell{가} & \makecell{1} & \makecell{102} & \makecell{대} & \makecell{12.3} & \makecell{연와조} & \makecell{16.1} & \makecell{2.2} & \makecell{공유} \\
\hline
\makecell{3} & \makecell{1} & \makecell{D} & \makecell{가} & \makecell{2} & \makecell{201} & \makecell{대} & \makecell{24.6} & \makecell{연와조} & \makecell{32.2} & \makecell{4.4} &  \\
\hline
\makecell{4} & \makecell{1} & \makecell{E} & \makecell{가} & \makecell{2} & \makecell{202} & \makecell{대} & \makecell{24.6} & \makecell{연와조} & \makecell{32.2} & \makecell{4.4} &  \\
\hline
\makecell{5} & \makecell{1} & \makecell{F} & \makecell{가} & \makecell{3} & \makecell{301} & \makecell{대} & \makecell{24.6} & \makecell{연와조} & \makecell{32.2} & \makecell{4.4} &  \\
\hline
\makecell{6} & \makecell{1} & \makecell{G} & \makecell{가} & \makecell{3} & \makecell{302} & \makecell{대} & \makecell{12.3} & \makecell{연와조} & \makecell{16.1} & \makecell{2.2} & \makecell{공유} \\
\hline
\makecell{6} & \makecell{2} & \makecell{H} & \makecell{가} & \makecell{3} & \makecell{302} & \makecell{대} & \makecell{12.3} & \makecell{연와조} & \makecell{16.1} & \makecell{2.2} & \makecell{공유} \\
\hline
\makecell{7} & \makecell{1} & \makecell{I} & \makecell{나} & \makecell{1} & \makecell{101} & \makecell{대} & \makecell{25.2} & \makecell{연와조} & \makecell{32.8} & \makecell{4.4} &  \\
\hline
\makecell{8} & \makecell{1} & \makecell{J} & \makecell{나} & \makecell{1} & \makecell{102} & \makecell{대} & \makecell{12.6} & \makecell{연와조} & \makecell{16.4} & \makecell{2.2} & \makecell{공유} \\
\hline
\makecell{8} & \makecell{2} & \makecell{K} & \makecell{나} & \makecell{1} & \makecell{102} & \makecell{대} & \makecell{12.6} & \makecell{연와조} & \makecell{16.4} & \makecell{2.2} & \makecell{공유} \\
\hline
\makecell{9} & \makecell{1} & \makecell{L} & \makecell{나} & \makecell{2} & \makecell{201} & \makecell{대} & \makecell{25.2} & \makecell{연와조} & \makecell{32.8} & \makecell{4.4} &  \\
\hline
\makecell{10} & \makecell{1} & \makecell{M} & \makecell{나} & \makecell{2} & \makecell{202} & \makecell{대} & \makecell{25.2} & \makecell{연와조} & \makecell{32.8} & \makecell{4.4} &  \\
\hline
\makecell{11} & \makecell{1} & \makecell{N} & \makecell{나} & \makecell{3} & \makecell{301} & \makecell{대} & \makecell{25.2} & \makecell{연와조} & \makecell{32.8} & \makecell{4.4} &  \\
\hline
\makecell{12} & \makecell{1} & \makecell{O} & \makecell{나} & \makecell{3} & \makecell{302} & \makecell{대} & \makecell{25.2} & \makecell{연와조} & \makecell{32.8} & \makecell{4.4} &  \\
\hline
\end{tabular}}\par}

< 자료 4 > 거래사례 \par
1. 거래사례(1) \par
1) 소 재 지: 평가동 1-1번지 ○○빌라 가동 1층 102호(사용승인일: 1985.01.02) \par
2) 거래시점: 2020.03.10 \par
3) 거래금액: 300,000,000원 \par
4) 건물내역: 전유(32.2m²), 공용(4.4m²), 소유권대지권(24.6m²), 남향, 계단식 \par
5) 기타사항: 당해 정비사업에 따른 개발이익이 포함되지 않은 정상거래사례임 \par
2. 거래사례(2) \par
1) 소 재 지: 평가동 1-2번지 ○○빌라 나동 2층 202호(사용승인일: 1985.01.07) \par
2) 거래시점: 2022.10.10 \par
3) 거래금액: 450,000,000원 \par
4) 건물내역: 전유(32.8m²), 공용(4.4m²), 소유권대지권(25.2m²), 동향, 계단식 \par
5) 기타사항: 인테리어비용이 포함된 거래로 사업구역 내 가장 최근 거래사례임 \par
3. 거래사례(3) \par
1) 소 재 지: 평가동 6번지 일동빌라 2층 202호(사용승인일: 1982.04.20) \par
2) 거래시점: 2023.03.20 \par
3) 거래금액: 300,000,000원 \par
4) 건물내역: 전유(29.5m²), 공용(4.05m²), 소유권대지권(22.5m²), 서향, 복도식 \par
5) 기타사항: 2인이 공유로 매수한 정상거래사례임 \par
4. 거래사례(4) \par
1) 소 재 지: 평가동 5번지 이동빌라 1층 101호(사용승인일: 1988.02.10) \par
2) 거래시점: 2023.04.28 \par
3) 거래금액: 350,000,000원 \par
4) 건물내역: 전유(30m²), 공용(3.05m²), 소유권대지권(10m²), 남향, 계단식 \par
5) 기타사항: 사례는 △△가로주택정비사업을 위한 추진위원회가 구성되어 있음 \par
5. 거래사례(5) \par
1) 소 재 지: 평가동 7번지 삼동빌라 2층 204호(사용승인일: 1987.03.16) \par
2) 거래시점: 2023.07.10 \par
3) 거래금액: 500,000,000원 \par
4) 건물내역: 전유(33.5m²), 공용(5.5m²), 소유권대지권(25.5m²), 동향, 복도식 \par
5) 기타사항: 법인이 소유한 물건으로 정상거래사례임 \par
< 자료 5 > 시점수정 자료 \par
{\small\centering
\resizebox{\ifdim\width>\linewidth\linewidth\else\width\fi}{!}{%
\begin{tabular}{|c|c|c|c|c|}
\hline
\makecell{거래사례(1)} & \makecell{거래사례(2)} & \makecell{거래사례(3)} & \makecell{거래사례(4)} & \makecell{거래사례(5)} \\
\hline
\makecell{1.05131} & \makecell{0.96353} & \makecell{0.99209} & \makecell{0.99702} & \makecell{0.98722} \\
\hline
\end{tabular}}\par}

< 자료 6 > 지역요인 자료 \par
평가대상과 거래사례는 인근지역에 소재하여 지역요인은 대등함 \par
< 자료 7 > 가치형성요인 자료 \par
1. 외부요인 \par
{\small\centering
\resizebox{\ifdim\width>\linewidth\linewidth\else\width\fi}{!}{%
\begin{tabular}{|c|c|c|c|c|}
\hline
\makecell{거래사례(1)} & \makecell{거래사례(2)} & \makecell{거래사례(3)} & \makecell{거래사례(4)} & \makecell{거래사례(5)} \\
\hline
\makecell{1.00} & \makecell{1.00} & \makecell{1.05} & \makecell{1.10} & \makecell{0.95} \\
\hline
\end{tabular}}\par}

2. 내부요인 \par
{\small\centering
\resizebox{\ifdim\width>\linewidth\linewidth\else\width\fi}{!}{%
\begin{tabular}{|c|c|c|c|c|}
\hline
\makecell{거래사례(1)} & \makecell{거래사례(2)} & \makecell{거래사례(3)} & \makecell{거래사례(4)} & \makecell{거래사례(5)} \\
\hline
\makecell{1.00} & \makecell{1.00} & \makecell{0.97} & \makecell{1.00} & \makecell{1.02} \\
\hline
\end{tabular}}\par}

3. 호별요인 \par
1) 층별 효용비율 \par
{\small\centering
\resizebox{\ifdim\width>\linewidth\linewidth\else\width\fi}{!}{%
\begin{tabular}{|c|c|c|c|c|c|}
\hline
\makecell{구분} & \makecell{지하1층} & \makecell{1층} & \makecell{2층} & \makecell{3층} & \makecell{4층} \\
\hline
\makecell{지하1층} & \makecell{1.00} & \makecell{1.11} & \makecell{1.17} & \makecell{1.14} & \makecell{1.09} \\
\hline
\makecell{1층} & \makecell{0.90} & \makecell{1.00} & \makecell{1.05} & \makecell{1.03} & \makecell{0.98} \\
\hline
\makecell{2층} & \makecell{0.86} & \makecell{0.95} & \makecell{1.00} & \makecell{0.98} & \makecell{0.93} \\
\hline
\makecell{3층} & \makecell{0.87} & \makecell{0.97} & \makecell{1.02} & \makecell{1.00} & \makecell{0.95} \\
\hline
\makecell{4층} & \makecell{0.92} & \makecell{1.02} & \makecell{1.07} & \makecell{1.05} & \makecell{1.00} \\
\hline
\end{tabular}}\par}

2) 향별 효용비율 \par
{\small\centering
\resizebox{\ifdim\width>\linewidth\linewidth\else\width\fi}{!}{%
\begin{tabular}{|c|c|c|c|c|}
\hline
\makecell{구분} & \makecell{동향} & \makecell{서향} & \makecell{남향} & \makecell{북향} \\
\hline
\makecell{동향} & \makecell{1.00} & \makecell{0.98} & \makecell{1.02} & \makecell{0.95} \\
\hline
\makecell{서향} & \makecell{1.02} & \makecell{1.00} & \makecell{1.04} & \makecell{0.97} \\
\hline
\makecell{남향} & \makecell{0.98} & \makecell{0.96} & \makecell{1.00} & \makecell{0.93} \\
\hline
\makecell{북향} & \makecell{1.05} & \makecell{1.03} & \makecell{1.07} & \makecell{1.00} \\
\hline
\end{tabular}}\par}

3) 형태별 효용비율 \par
{\small\centering
\resizebox{\ifdim\width>\linewidth\linewidth\else\width\fi}{!}{%
\begin{tabular}{|c|c|c|}
\hline
\makecell{구분} & \makecell{계단식} & \makecell{복도식} \\
\hline
\makecell{계단식} & \makecell{1.00} & \makecell{0.95} \\
\hline
\makecell{복도식} & \makecell{1.05} & \makecell{1.00} \\
\hline
\end{tabular}}\par}

4) 주거환경 영향지수 \par
주거환경 영향지수는 전문기관이 수행한 다음 자료를 적용 \par
{\small\centering
\resizebox{\ifdim\width>\linewidth\linewidth\else\width\fi}{!}{%
\begin{tabular}{|c|c|c|c|c|c|}
\hline
\makecell{구분} & \makecell{일조\\지수} & \makecell{조망\\지수} & \makecell{소음\\지수} & \makecell{사생활\\침해지수} & \makecell{주거환경\\영향지수} \\
\hline
\makecell{명부 연번 1} & \makecell{0.97} & \makecell{0.95} & \makecell{0.99} & \makecell{1.01} & \makecell{0.92} \\
\hline
\makecell{명부 연번 2} & \makecell{0.98} & \makecell{0.95} & \makecell{0.99} & \makecell{1.01} & \makecell{0.93} \\
\hline
\makecell{명부 연번 3} & \makecell{0.98} & \makecell{0.96} & \makecell{1.00} & \makecell{1.01} & \makecell{0.95} \\
\hline
\makecell{명부 연번 4} & \makecell{0.98} & \makecell{0.96} & \makecell{1.01} & \makecell{1.01} & \makecell{0.96} \\
\hline
\makecell{명부 연번 5} & \makecell{0.98} & \makecell{0.97} & \makecell{1.01} & \makecell{1.02} & \makecell{0.98} \\
\hline
\makecell{명부 연번 6} & \makecell{0.99} & \makecell{0.97} & \makecell{1.01} & \makecell{1.02} & \makecell{0.99} \\
\hline
\makecell{명부 연번 7} & \makecell{0.98} & \makecell{0.97} & \makecell{0.98} & \makecell{1.01} & \makecell{0.94} \\
\hline
\makecell{명부 연번 8} & \makecell{0.98} & \makecell{0.97} & \makecell{0.99} & \makecell{1.01} & \makecell{0.95} \\
\hline
\makecell{명부 연번 9} & \makecell{0.99} & \makecell{0.97} & \makecell{0.99} & \makecell{1.02} & \makecell{0.97} \\
\hline
\makecell{명부 연번 10} & \makecell{0.99} & \makecell{1.00} & \makecell{0.99} & \makecell{1.02} & \makecell{1.00} \\
\hline
\makecell{명부 연번 11} & \makecell{1.01} & \makecell{1.00} & \makecell{0.99} & \makecell{1.02} & \makecell{1.02} \\
\hline
\makecell{명부 연번 12} & \makecell{1.01} & \makecell{1.00} & \makecell{0.99} & \makecell{1.05} & \makecell{1.05} \\
\hline
\makecell{거래사례(3)} & \makecell{0.98} & \makecell{0.99} & \makecell{1.02} & \makecell{0.98} & \makecell{0.97} \\
\hline
\makecell{거래사례(4)} & \makecell{0.97} & \makecell{0.97} & \makecell{0.99} & \makecell{1.01} & \makecell{0.94} \\
\hline
\makecell{거래사례(5)} & \makecell{0.99} & \makecell{1.01} & \makecell{0.99} & \makecell{0.99} & \makecell{0.98} \\
\hline
\end{tabular}}\par}

4. 기타요인 \par
제시된 자료 이외 기타 가격에 영향을 미치는 요인은 없음 \par
< 자료 8 > 종후자산(아파트) 감정평가액 \par
{\small\centering
\resizebox{\ifdim\width>\linewidth\linewidth\else\width\fi}{!}{%
\begin{tabular}{|c|c|c|c|c|c|c|c|c|}
\hline
\makecell{구분} & \makecell{동} & \makecell{층} & \makecell{호} & \makecell{타입} & \makecell{전유\\(m²)} & \makecell{공용\\(m²)} & \makecell{공급\\(m²)} & \makecell{감정평가액\\(원)} \\
\hline
\makecell{1} & \makecell{101} & \makecell{2} & \makecell{201} & \makecell{40A} & \makecell{40.11} & \makecell{12.3} & \makecell{52.41} & \makecell{422,000,000} \\
\hline
\makecell{2} & \makecell{101} & \makecell{3} & \makecell{301} & \makecell{40A} & \makecell{40.11} & \makecell{12.3} & \makecell{52.41} & \makecell{426,000,000} \\
\hline
\makecell{3} & \makecell{101} & \makecell{4} & \makecell{401} & \makecell{40A} & \makecell{40.11} & \makecell{12.3} & \makecell{52.41} & \makecell{439,000,000} \\
\hline
\makecell{4} & \makecell{101} & \makecell{5} & \makecell{501} & \makecell{40A} & \makecell{40.11} & \makecell{12.3} & \makecell{52.41} & \makecell{454,000,000} \\
\hline
\makecell{5} & \makecell{101} & \makecell{6} & \makecell{601} & \makecell{40A} & \makecell{40.11} & \makecell{12.3} & \makecell{52.41} & \makecell{472,000,000} \\
\hline
\makecell{6} & \makecell{101} & \makecell{7} & \makecell{701} & \makecell{40A} & \makecell{40.11} & \makecell{12.3} & \makecell{52.41} & \makecell{491,000,000} \\
\hline
\makecell{7} & \makecell{101} & \makecell{8} & \makecell{801} & \makecell{40A} & \makecell{40.11} & \makecell{12.3} & \makecell{52.41} & \makecell{516,000,000} \\
\hline
\makecell{8} & \makecell{101} & \makecell{9} & \makecell{901} & \makecell{40A} & \makecell{40.11} & \makecell{12.3} & \makecell{52.41} & \makecell{539,000,000} \\
\hline
\makecell{9} & \makecell{101} & \makecell{2} & \makecell{202} & \makecell{40B} & \makecell{40.11} & \makecell{12.3} & \makecell{52.41} & \makecell{425,000,000} \\
\hline
\makecell{10} & \makecell{101} & \makecell{3} & \makecell{302} & \makecell{40B} & \makecell{40.11} & \makecell{12.3} & \makecell{52.41} & \makecell{429,000,000} \\
\hline
\makecell{11} & \makecell{101} & \makecell{4} & \makecell{402} & \makecell{40B} & \makecell{40.11} & \makecell{12.3} & \makecell{52.41} & \makecell{442,000,000} \\
\hline
\makecell{12} & \makecell{101} & \makecell{5} & \makecell{502} & \makecell{40B} & \makecell{40.11} & \makecell{12.3} & \makecell{52.41} & \makecell{457,000,000} \\
\hline
\makecell{13} & \makecell{101} & \makecell{6} & \makecell{602} & \makecell{40B} & \makecell{40.11} & \makecell{12.3} & \makecell{52.41} & \makecell{476,000,000} \\
\hline
\makecell{14} & \makecell{101} & \makecell{7} & \makecell{702} & \makecell{40B} & \makecell{40.11} & \makecell{12.3} & \makecell{52.41} & \makecell{495,000,000} \\
\hline
\makecell{15} & \makecell{101} & \makecell{8} & \makecell{802} & \makecell{40B} & \makecell{40.11} & \makecell{12.3} & \makecell{52.41} & \makecell{519,000,000} \\
\hline
\makecell{16} & \makecell{101} & \makecell{9} & \makecell{902} & \makecell{40B} & \makecell{40.11} & \makecell{12.3} & \makecell{52.41} & \makecell{543,000,000} \\
\hline
\makecell{17} & \makecell{101} & \makecell{2} & \makecell{203} & \makecell{40C} & \makecell{40.11} & \makecell{12.3} & \makecell{52.41} & \makecell{424,000,000} \\
\hline
\makecell{18} & \makecell{101} & \makecell{3} & \makecell{303} & \makecell{40C} & \makecell{40.11} & \makecell{12.3} & \makecell{52.41} & \makecell{428,000,000} \\
\hline
\makecell{19} & \makecell{101} & \makecell{4} & \makecell{403} & \makecell{40C} & \makecell{40.11} & \makecell{12.3} & \makecell{52.41} & \makecell{441,000,000} \\
\hline
\makecell{20} & \makecell{101} & \makecell{5} & \makecell{503} & \makecell{40C} & \makecell{40.11} & \makecell{12.3} & \makecell{52.41} & \makecell{456,000,000} \\
\hline
\makecell{21} & \makecell{101} & \makecell{6} & \makecell{603} & \makecell{40C} & \makecell{40.11} & \makecell{12.3} & \makecell{52.41} & \makecell{475,000,000} \\
\hline
\makecell{22} & \makecell{101} & \makecell{7} & \makecell{703} & \makecell{40C} & \makecell{40.11} & \makecell{12.3} & \makecell{52.41} & \makecell{494,000,000} \\
\hline
\makecell{23} & \makecell{101} & \makecell{8} & \makecell{803} & \makecell{40C} & \makecell{40.11} & \makecell{12.3} & \makecell{52.41} & \makecell{518,000,000} \\
\hline
\makecell{24} & \makecell{101} & \makecell{9} & \makecell{903} & \makecell{40C} & \makecell{40.11} & \makecell{12.3} & \makecell{52.41} & \makecell{542,000,000} \\
\hline
\multicolumn{5}{|c|}{\makecell{합계}} & \makecell{962.64} & \makecell{295.2} & \makecell{1,257.84} & \makecell{11,323,000,000} \\
\hline
\end{tabular}}\par}

< 자료 9 > 정비사업비 관련 제시자료 \par
{\small\centering
\resizebox{\ifdim\width>\linewidth\linewidth\else\width\fi}{!}{%
\begin{tabular}{|c|c|c|c|}
\hline
\multicolumn{2}{|c|}{\makecell{항목}} & \makecell{금액(원)} & \makecell{귀속} \\
\hline
\multirow{4}{*}{\makecell{토지}} & \makecell{원가산입 종전자산가액} & \makecell{물음 1) 적용} & \makecell{토지} \\
\cline{2-4}
 & \makecell{시유지 매입비} & \makecell{210,000,000} & \makecell{토지} \\
\cline{2-4}
 & \makecell{측량 및 지질조사 등} & \makecell{110,000,000} & \makecell{토지} \\
\cline{2-4}
 & \makecell{취득세 및 등록세 등} & \makecell{155,000,000} & \makecell{토지} \\
\hline
\multirow{3}{*}{\makecell{건축비}} & \makecell{직접공사비} & \makecell{4,520,000,000} & \makecell{건물} \\
\cline{2-4}
 & \makecell{간접공사비(인입비, 부담금 등)} & \makecell{25,000,000} & \makecell{건물} \\
\cline{2-4}
 & \makecell{설계비, 감리비 등} & \makecell{300,000,000} & \makecell{건물} \\
\hline
\multirow{3}{*}{\makecell{부대경비}} & \makecell{외주용역비(정비사업, 감정평가 등)} & \makecell{150,000,000} & \makecell{공통} \\
\cline{2-4}
 & \makecell{회계감사비, 신탁등기비, 보존등기비 등} & \makecell{265,000,000} & \makecell{공통} \\
\cline{2-4}
 & \makecell{이주정착금, 예비비 등} & \makecell{500,000,000} & \makecell{공통} \\
\hline
\multirow{3}{*}{\makecell{금융비용}} & \makecell{대출수수료} & \makecell{150,000,000} & \makecell{공통} \\
\cline{2-4}
 & \makecell{조합원이주이자} & \makecell{100,000,000} & \makecell{공통} \\
\cline{2-4}
 & \makecell{사업비이자 등} & \makecell{150,000,000} & \makecell{공통} \\
\hline
\multicolumn{2}{|c|}{\makecell{합계(원가산입 종전자산가액 제외)}} & \makecell{6,635,000,000} & \makecell{-} \\
\hline
\end{tabular}}\par}

\end{longquestion}
